%%%%%%%%%%%%%%%%%%%%%%%%%%%%%%%%%%%%%%%%%%%%%%%%%%%%%
%                                                   %
% Cette version repart de 03_Draft\conclusionV1.tex %
%                                                   %
%%%%%%%%%%%%%%%%%%%%%%%%%%%%%%%%%%%%%%%%%%%%%%%%%%%%%

% NB: Je supprime toutes les (grosses) parties qui étaient en commentaire dans la version précédente pour qu'on y voit plus clair.

In this paper, we show that displaced workers in France experience large and persistent earnings losses. Lower hourly wage rates after re-employment account for a larger share of earnings losses over time. In turn, losses in employer-specific wage premia account for a substantial share of the decline in hourly wage rates. Contrasting these losses in wage premia with measures of productivity, we show that displaced workers tend to be reemployed by more productive firms with a lower labor share of value added. These findings suggest that losses in firm wage premia are not driven by reallocation toward low-surplus firms, but rather by reemployment in firms with less favorable wage policies given their productivity. This is suggestive of a loss in negotiation power which is confirmed by direct evidence on qualitative variables capturing the quality of labor relations at the firm-level.  

Overall, our results show that job displacements can contribute to reallocate workers to more productive firms but that this process is costly for workers, in part due to the fact that destination firms have across-the-board less favorable wage policy. From a positive and normative standpoint, understanding the origin of these firms' ability to constrain labor compensation would be important.\footnote{This could derive from greater monopsonistic power which in turn can be measured by firm-specific markdowns in the spirit of \cite{hershbein_macaluso_monopsony_2020}. However, this methodology is best suited to the manufacturing sector, which is an issue in our context where a large fraction (60\% see Figure \ref{fig:desc:LEquali}) of our sample transitions inside and outside the manufacturing sector.} On the policy side, as reallocations seem to provide a social good, it would be useful to explore ways of mitigating the private cost they impose on workers, for instance through wage insurance. %as an alternative to standard unemployment insurance systems
While beyond the scope of this paper, we view both questions as interesting avenues for further research.

